% GENERATED FILE. Edit canonical lessons, then run npm run generate:books.
% canonical-source-hash: fnv1a64:48c6af8b4af59a02
% canonical-lessons: MR-C532-buddhibal, MR-C532-songti, MR-C532-phasa, MR-C532-sabdakode, MR-C532-lapachpi

\chapter{Games at Home}
\label{ch:mr-games-at-home}
\hlchaptermodality{532}

\begin{chapteropening}
\textbf{By the end of this chapter:} \emph{I can talk about board games, dice and puzzles.}

Complete the last lesson of chapter 532: I can talk about board games, dice and puzzles.
\end{chapteropening}

\section[\texorpdfstring{\mr{बुद्धिबळ} (buddhibaḷ)}{बुद्धिबळ (buddhibaḷ)}]{\mr{बुद्धिबळ} (buddhibaḷ) — chess}

\label{lesson:MR-C532-buddhibal}



\begin{quote}
Before the new one: say the Marathi for to hit away, then the Marathi for to catch a ball.
\end{quote}

\subsection*{You'll want to know: \mr{बुद्धिबळ}}
\textbf{\mr{बुद्धिबळ}} — \emph{buddhibaḷ} — \textquotedblleft{}chess\textquotedblright{}.

\begin{quote}
\mr{आजोबा} \mr{आणि} \mr{मी} \mr{रविवारी} \mr{बुद्धिबळ} \mr{खेळतो}.
\end{quote}

\begin{quote}
\emph{ājobā āṇi mī ravivārī buddhibaḷ kheḷto.}
\end{quote}

\begin{quote}
Grandfather and I play chess on Sundays.
\end{quote}

Neuter.

\begin{cousinweb}
\textbf{\mr{बुद्धिबळ}} (\emph{buddhibaḷ}) is a tatsama, the Sanskrit compound \emph{buddhibala}, \emph{buddhi}, \textquotedblleft{}intellect\textquotedblright{}, and \emph{bala}, \textquotedblleft{}strength\textquotedblright{}, listed as the name of a kind of game (Monier-Williams); Marathi uses it for chess.
\end{cousinweb}

\begin{grammarlens}[title={What you've built}]
One more word for describing what you lived and felt.
\end{grammarlens}

\subsection*{Your turn}
\begin{itemize}
\raggedright
  \item \emph{Say it:} \emph{buddhibaḷ}
  \item \emph{Say it:} \emph{buddhibaḷ}, once more
  \item \emph{Say it:} \emph{ājobā āṇi mī ravivārī buddhibaḷ kheḷto.} — Grandfather and I play chess on Sundays.
  \item \emph{Recall:} say the Marathi for to hit away, then the Marathi for to catch a ball, then say \emph{buddhibaḷ} again
\end{itemize}

\begin{tcolorbox}[breakable,colback=teal!4,colframe=teal!35!black,title={Before you move on}]
What does \mr{बुद्धिबळ} mean? (\textquotedblleft{}Chess\textquotedblright{}.) Say it once more.
\end{tcolorbox}
\section[\texorpdfstring{\mr{सोंगटी} (soṅgṭī)}{सोंगटी (soṅgṭī)}]{\mr{सोंगटी} (soṅgṭī) — a game piece; a counter}

\label{lesson:MR-C532-songti}



\begin{quote}
Before the new one: say the Marathi for to catch a ball, then the Marathi for chess.
\end{quote}

\subsection*{You'll want to know: \mr{सोंगटी}}
\textbf{\mr{सोंगटी}} — \emph{soṅgṭī} — \textquotedblleft{}a game piece; a counter\textquotedblright{}.

\begin{quote}
\mr{बुद्धिबळाची} \mr{एक} \mr{सोंगटी} \mr{हरवली}.
\end{quote}

\begin{quote}
\emph{buddhibaḷācī ek soṅgṭī haravlī.}
\end{quote}

\begin{quote}
One of the chess pieces got lost.
\end{quote}

Feminine.

\begin{cousinweb}
\textbf{\mr{सोंगटी}} (\emph{soṅgṭī}) is the word for a piece in a board game; its origin is not settled.
\end{cousinweb}

\begin{grammarlens}[title={What you've built}]
One more word for describing what you lived and felt.
\end{grammarlens}

\subsection*{Your turn}
\begin{itemize}
\raggedright
  \item \emph{Say it:} \emph{soṅgṭī}
  \item \emph{Say it:} \emph{soṅgṭī}, once more
  \item \emph{Say it:} \emph{buddhibaḷācī ek soṅgṭī haravlī.} — One of the chess pieces got lost.
  \item \emph{Recall:} say the Marathi for to catch a ball, then the Marathi for chess, then say \emph{soṅgṭī} again
\end{itemize}

\begin{tcolorbox}[breakable,colback=teal!4,colframe=teal!35!black,title={Before you move on}]
What does \mr{सोंगटी} mean? (\textquotedblleft{}A game piece; a counter\textquotedblright{}.) Say it once more.
\end{tcolorbox}
\section[\texorpdfstring{\mr{फासा} (phāsā)}{फासा (phāsā)}]{\mr{फासा} (phāsā) — a die for games}

\label{lesson:MR-C532-phasa}



\begin{quote}
Before the new one: say the Marathi for chess, then the Marathi for a game piece.
\end{quote}

\subsection*{You'll want to know: \mr{फासा}}
\textbf{\mr{फासा}} — \emph{phāsā} — \textquotedblleft{}a die for games\textquotedblright{}.

\begin{quote}
\mr{मी} \mr{फासा} \mr{टाकला} \mr{आणि} \mr{सहा} \mr{पडले}.
\end{quote}

\begin{quote}
\emph{mī phāsā ṭāklā āṇi sahā paḍle.}
\end{quote}

\begin{quote}
I threw the die and got a six.
\end{quote}

Masculine. Plural: \mr{फासे}, \emph{phāse}.

\mr{फासा} \mr{टाकणे}, to throw the die.

\begin{cousinweb}
\textbf{\mr{फासा}} (\emph{phāsā}) is a tadbhava: CDIAL lists Marathi \emph{phāsā}, \textquotedblleft{}die\textquotedblright{}, under Sanskrit \emph{pāśa}, \textquotedblleft{}die, dice\textquotedblright{} (CDIAL 8132).
\end{cousinweb}

\begin{grammarlens}[title={What you've built}]
One more word for describing what you lived and felt.
\end{grammarlens}

\subsection*{Your turn}
\begin{itemize}
\raggedright
  \item \emph{Say it:} \emph{phāsā}
  \item \emph{Say it:} \emph{phāsā}, once more
  \item \emph{Say it:} \emph{mī phāsā ṭāklā āṇi sahā paḍle.} — I threw the die and got a six.
  \item \emph{Recall:} say the Marathi for chess, then the Marathi for a game piece, then say \emph{phāsā} again
\end{itemize}

\begin{tcolorbox}[breakable,colback=teal!4,colframe=teal!35!black,title={Before you move on}]
What does \mr{फासा} mean? (\textquotedblleft{}A die for games\textquotedblright{}.) Say it once more.
\end{tcolorbox}
\section[\texorpdfstring{\mr{शब्दकोडे} (śabdakoḍe)}{शब्दकोडे (śabdakoḍe)}]{\mr{शब्दकोडे} (śabdakoḍe) — a crossword}

\label{lesson:MR-C532-sabdakode}



\begin{quote}
Before the new one: say the Marathi for a game piece, then the Marathi for a die for games.
\end{quote}

\subsection*{You'll want to know: \mr{शब्दकोडे}}
\textbf{\mr{शब्दकोडे}} — \emph{śabdakoḍe} — \textquotedblleft{}a crossword\textquotedblright{}.

\begin{quote}
\mr{आजोबा} \mr{रोज} \mr{सकाळी} \mr{एक} \mr{शब्दकोडं} \mr{सोडवतात}.
\end{quote}

\begin{quote}
\emph{ājobā roj sakāḷī ek śabdakoḍaṁ soḍavtāt.}
\end{quote}

\begin{quote}
Grandfather solves a crossword every morning.
\end{quote}

Neuter.

\mr{शब्दकोडं} \mr{सोडवणे}, to solve a crossword; the dictionary writes \mr{शब्दकोडे}.

\begin{cousinweb}
\textbf{\mr{शब्दकोडे}} (\emph{śabdakoḍe}) is a Marathi compound built from \mr{शब्द}, \textquotedblleft{}word\textquotedblright{}, from Sanskrit \emph{śabda} (Monier-Williams), and \mr{कोडे}, \textquotedblleft{}a puzzle\textquotedblright{}.
\end{cousinweb}

\begin{grammarlens}[title={What you've built}]
One more word for describing what you lived and felt.
\end{grammarlens}

\subsection*{Your turn}
\begin{itemize}
\raggedright
  \item \emph{Say it:} \emph{śabdakoḍe}
  \item \emph{Say it:} \emph{śabdakoḍe}, once more
  \item \emph{Say it:} \emph{ājobā roj sakāḷī ek śabdakoḍaṁ soḍavtāt.} — Grandfather solves a crossword every morning.
  \item \emph{Recall:} say the Marathi for a game piece, then the Marathi for a die for games, then say \emph{śabdakoḍe} again
\end{itemize}

\begin{tcolorbox}[breakable,colback=teal!4,colframe=teal!35!black,title={Before you move on}]
What does \mr{शब्दकोडे} mean? (\textquotedblleft{}A crossword\textquotedblright{}.) Say it once more.
\end{tcolorbox}
\section[\texorpdfstring{\mr{लपाछपी} (lapāchpī)}{लपाछपी (lapāchpī)}]{\mr{लपाछपी} (lapāchpī) — hide and seek}

\label{lesson:MR-C532-lapachpi}



\begin{quote}
Before the new one: say the Marathi for a die for games, then the Marathi for a crossword.
\end{quote}

\subsection*{You'll want to know: \mr{लपाछपी}}
\textbf{\mr{लपाछपी}} — \emph{lapāchpī} — \textquotedblleft{}hide and seek\textquotedblright{}.

\begin{quote}
\mr{सुट्टीत} \mr{मुलं} \mr{अंगणात} \mr{लपाछपी} \mr{खेळली}.
\end{quote}

\begin{quote}
\emph{suṭṭīt mulaṁ aṅgaṇāt lapāchpī kheḷlī.}
\end{quote}

\begin{quote}
In the holidays the children played hide and seek in the yard.
\end{quote}

Feminine.

\begin{cousinweb}
\textbf{\mr{लपाछपी}} (\emph{lapāchpī}) is a rhyming Marathi compound built on \mr{लपणे}, \textquotedblleft{}to hide\textquotedblright{}; the origin of the second half is not settled.
\end{cousinweb}

\begin{grammarlens}[title={What you've built}]
One more word for describing what you lived and felt.
\end{grammarlens}

\subsection*{Your turn}
\begin{itemize}
\raggedright
  \item \emph{Say it:} \emph{lapāchpī}
  \item \emph{Say it:} \emph{lapāchpī}, once more
  \item \emph{Say it:} \emph{suṭṭīt mulaṁ aṅgaṇāt lapāchpī kheḷlī.} — In the holidays the children played hide and seek in the yard.
  \item \emph{Recall:} say the Marathi for a die for games, then the Marathi for a crossword, then say \emph{lapāchpī} again
\end{itemize}

\begin{tcolorbox}[breakable,colback=teal!4,colframe=teal!35!black,title={Before you move on}]
What does \mr{लपाछपी} mean? (\textquotedblleft{}Hide and seek\textquotedblright{}.) Say it once more.
\end{tcolorbox}
